\section{Methodology}

Building on our observation that permutation invariance must be architecturally encoded, we design experiments to verify both the existence of equivariance benefits and the mechanism underlying them.

\subsection{Neural Functional Networks (NFN)}

We adopt the NFN architecture from Zhou et al.~\cite{zhou2023neural}. The core building block is NPLinear, which linearly transforms weight tensors while maintaining equivariance. HNPPool produces permutation-invariant representations via pooling over neuron dimensions.

Our NFNRegressor:
\begin{verbatim}
Input weights -> NPLinear(32) -> ReLU
  -> NPLinear(32) -> ReLU -> HNPPool -> Linear(1)
\end{verbatim}

\subsection{MLP-Matched Baseline}

For fair comparison, we design an MLP with matched capacity:
\begin{verbatim}
Flattened weights -> Linear(256) -> ReLU
  -> Linear(256) -> ReLU -> Linear(1)
\end{verbatim}

\subsection{Probe Invariance Test}

To verify whether models develop permutation-invariant representations:

\begin{enumerate}
    \item Take a trained model $M$ and a test weight tensor $W$
    \item Generate $K=10$ random permutations $\pi_1, \pi_2, \ldots, \pi_K$
    \item Apply $M$ to each permuted version: $\hat{y}_i = M(\pi_i(W))$
    \item Compute invariance score: $I = \text{mean pairwise correlation}$
\end{enumerate}

For mathematically invariant models: $I = 1.0$.
